\newpage


\section{Введение}
Решается задача выбора оптимальной структуры нейронной сети. В силу высокой вычислительной сложности, время оптимизации нейронных сетей может занимать до нескольких дней. Поэтому построение и выбороптимальной структуры нейронной сети также является вычислительно сложной процедурой, которая значимо влияет на итоговое качество
модели. Использование избыточно сложных моделей с избыточным числом неинформативных параметров является препятствием для использования глубоких сетей на мобильных устройствах в режиме реального времени. Существует ряд подходов к построению оптимальной сети. В работах [2, 3] предлагается использовать модель градиентного спуска для оптимизации сети. В [4] используются байесовские методы оптимизации параметров нейронных сетей. Другим методом поиска оптимальной структуры является прореживание избыточно сложной модели.
В работе [6] предлагается удалять наименее релевантные параметры на основе значений первой и второй производных функции ошибки.
Данная работа посвящена прореживанию структуры сети. Предлагается удалять наименее релевантные параметры модели. Под релевантностью [6] подразумевается то, насколько параметр влияет на функцию
ошибки. Малая релевантность указывает на то, что удаление этого параметра не влечет значимого изменения функции ошибки. Метод предлагает построение исходной избыточной сложности нейросети с большим числом избыточных параметров. Для определения релевантности параметров предлагается оптимизацировать параметры и гиперпараметры в
единой процедуре. Для удаления параметров предлагается использовать метод Белсли.
Проверка и анализ метода проводится на выборке Boston Housing, Wine и синтетических данных. Результат сравнивается с моделью, полученной при помощи базовых алгоритмов.

\section{Анализ литературы}
To be done...
